% Abstract -- front matter (\input by main.tex before Chapter 1)
% Metric-first funnel (four short paragraphs): the measurement gap (Chamfer invalid
% on real LiDAR, #26) -> the contribution (the donor metric + its four-part
% validation + the completion payoff, PARTIALLY-HOLDS held-out) -> the pipeline it
% lives in and the two diagnostic findings (recall root-cause; synthetic-pretraining
% redundancy) with the segmentation bottleneck as one late sentence -> the
% three-contribution close. It is the first thing a committee reads, so every number
% and hedge must be defensible. It cites no chapters and no literature -- an abstract
% is self-contained prose -- so there are no cross-references or \cite{} here.
%
% Honesty-phrasing rules applied (personal_agenda P1):
%   - recall (0.730) carries its qualifier inline ("at a point-IoU of 0.3, on cars
%     with enough returns to be scorable") so it is NOT read as recall against all
%     annotated cars (~0.54, Finding #48); the exact >=10-surviving-point denominator
%     mechanics live in Sec 4.1
%   - LOSO is stated qualitatively ("validated across all eleven labelled sequences to
%     rule out model-selection bias"); the cross-validated numbers (0.737/0.719) live in
%     Sec 4.1 -- abstract-level compression, not a softening
%   - the amodal-box benefit is scoped to "static, gate-passed cars" (movers are
%     plausibility-only #44; long band untested off seq 08)
%   - the held-out completion result is "partially replicates", never "confirmed"/"holds"
%   - the detection headline is a single-sequence (seq 08) result, cross-validated
%   - negatives (recall repairs, synthetic redundancy) are named as findings
%   - the donor metric is "a validated" metric, NOT "the first"
%     (no literature-priority overclaim; matches C7 and the Ch 9 review)
%   - the three closing contributions match Ch 1 Sec 1.4 / Ch 9 Sec 9.3 / Ch 2 Sec 2.5
%     EXACTLY (donor metric / recall root-cause / synthetic-pretraining redundancy);
%     claims kept within the Section-6 claims map (C1-C11); no mover accuracy claim
%
\chapter*{Abstract}
\label{ch:abstract}

Automotive LiDAR sees every car from one side, and point-cloud completion aims to
restore the occluded surface --- but on real data it cannot be measured. No clean
reference shape exists for a car, and Chamfer distance against an accumulated
multi-frame cloud is invalid: the accumulation is occluded on the same side, so the raw
partial scores \emph{best}, ranking a network that does nothing above one that completes.

The central contribution of this thesis is a valid metric: a donor-frame coverage score
for how much of the surface seen only in \emph{other} frames of the same static car a
completion recovers, guarded against out-of-box hallucination and compared to a
parameter-free symmetry-mirror baseline, and passing a four-part validation battery. On
real sequence-08 cars a PCN network recovers genuinely unseen surface well above the
mirror and tightens amodal boxes on static, gate-passed cars --- a benefit that
partially replicates on a pre-registered held-out sequence.

The metric lives in an interpretable, near-annotation-free pipeline of classical
geometry and small learned models. Characterising it yields two diagnostic findings:
detection is high-precision (\num{0.905}) but structurally recall-limited (\num{0.730}
at a point-IoU of \num{0.3}, on cars with enough returns to be scorable; cross-validated
across all eleven labelled sequences), root-caused to clustering fragmenting single
vehicles rather than the classifier; and synthetic pretraining is redundant once real
clusters are plentiful. The contributions are methodological and diagnostic, not
architectural.
